\documentclass[conference]{IEEEtran}

\usepackage{graphicx}
\usepackage{amsmath}
\usepackage{amssymb}
\usepackage{booktabs}
\usepackage{cite}
\usepackage{url}

\hyphenation{op-tical net-works case-based re-triev-al aug-men-ta-tion}

\def\casefig{../notebooks/paper/fig_case_distribution.png}
\def\compfig{../notebooks/paper/fig_retrieval_comparison.png}
\def\stabfig{../notebooks/paper/fig_retrieval_stability.png}

\begin{document}

\title{Experience-Driven Retrieval-Augmented Generation for Logistics Disruption Management}

\author{\IEEEauthorblockN{Ken}
\IEEEauthorblockA{School of Management\\
ken@example.edu}}

\maketitle

\begin{abstract}
Logistics disruption management requires decisions that are simultaneously fast, explainable, and robust to diverse operational contexts. Classical rule-based dispatching and similarity-based case retrieval remain attractive because they are inexpensive, but they often fail to capture semantic relations between heterogeneous disruption descriptions and therefore cannot reliably reuse prior operational experience. This paper presents an experience-driven retrieval-augmented generation framework for logistics disruption management, where a curated case memory is combined with dense retrieval, action-oriented reranking, and reward-aware reasoning. The system is evaluated on a case library of 1,000 historical disruption cases covering 21 event types, four operational scenarios, and five candidate actions. We test on 420 benchmark queries and compare a five-dimensional scoring baseline, a first retrieval-plus-prediction variant, and the proposed V2 method. V2 achieves 99.762\% action accuracy and 99.685\% NDCG@5 with an average decision time of 8.415\,ms, outperforming the baseline and V1 on both decision quality and retrieval quality. Repeated-query analysis also shows strong retrieval stability, with 99.905\% Top-1 action consistency. The results indicate that experience-driven RAG can deliver near-perfect disruption decisions while preserving millisecond-level responsiveness suitable for real-time logistics operations.
\end{abstract}

\begin{IEEEkeywords}
retrieval-augmented generation, logistics disruption management, case-based reasoning, dense retrieval, decision support, supply chain resilience
\end{IEEEkeywords}

\section{Introduction}
\label{sec:introduction}

Logistics systems operate under permanent exposure to disruptions such as vehicle breakdowns, warehouse delays, congestion, weather events, demand shocks, and last-minute order changes. These events alter route feasibility, service commitments, and resource utilization on very short time scales. In practice, dispatchers must decide whether to reroute, delay, ignore, reassign, or rebalance capacity before full information is available. The central challenge is therefore not only to optimize under uncertainty, but to do so quickly enough that the decision remains operationally relevant.

The traditional response in logistics analytics has been to combine optimization, heuristics, and hand-crafted decision rules. Vehicle-routing and dynamic-routing methods provide strong normative baselines for structured planning problems \cite{toth2014vrp,pillac2013dvrp}. They are indispensable when constraints are explicit and the cost model is reliable. However, disruption management is often triggered by semi-structured event descriptions and partial context: a dispatcher may know that a bridge restriction, a demand surge, and a vehicle maintenance alert occur simultaneously, yet may not have enough time to launch a full replanning cycle. In these settings, experience reuse is attractive because it compresses prior operational knowledge into directly actionable precedents.

Case-based reasoning (CBR) formalizes exactly this intuition: solve a new problem by retrieving, adapting, and retaining solutions from similar historical cases \cite{aamodt1994cbr,watson1999cbr}. In supply-chain settings, CBR has been used for collaborative planning, negotiation, and risk-related decisions \cite{kwon2007mace,fang2010hybrid,giannakis2011risk}. The weakness of classical CBR is that retrieval quality depends heavily on manually engineered similarity functions. This is problematic in disruption management because semantically related events may look lexically different, while superficially similar incidents can demand opposite actions once scale, severity, or scenario pressure is considered.

Recent progress in retrieval-augmented generation (RAG) offers a principled way to modernize experience reuse. RAG combines non-parametric memory with language-model reasoning so that retrieval supplies grounded evidence and generation provides contextual decision synthesis \cite{lewis2020rag}. Dense retrieval models further improve this pipeline by mapping queries and candidates into semantic vector spaces, reducing dependence on exact lexical overlap \cite{karpukhin2020dpr,reimers2019sbert,khattab2020colbert,santhanam2022colbertv2}. Meanwhile, lightweight large-language-model deployment and quantization methods make it increasingly practical to introduce reasoning components into latency-sensitive systems \cite{touvron2023llama,dettmers2023qlora,lin2024awq}. For logistics, this suggests a new design pattern: keep the decision anchored in historical operational memory, but use retrieval and generation to infer the best action under sparse and heterogeneous evidence.

This paper investigates that pattern for logistics disruption management. We study an experience-driven RAG system over a case library of 1,000 disruptions spanning 21 event types and four scenarios: small, medium, large, and stress. The library is operationally balanced at the scenario level, with 250 cases per scenario, but highly skewed by event type, where \textit{vehicle\_breakdown} alone accounts for 361 cases. Such imbalance mirrors practical operations: some disruptions dominate daily dispatching, while others are rare but high impact. The library also records a 76.7\% success rate, allowing the system to reason not only about what action was taken historically, but also whether that action succeeded.

We compare three decision mechanisms. The baseline uses a five-dimensional hand-crafted score. V1 augments retrieval with action prediction. V2, our final system, adds reward-aware reasoning over retrieved experience. Across 420 benchmark queries, V2 reaches 99.762\% action accuracy and 99.685\% NDCG@5 with 8.415\,ms average decision time. The key point is not merely that V2 is more accurate; it is that retrieval quality, action quality, and run-time all remain compatible with real-time operations.

The paper makes four contributions. First, it formulates logistics disruption management as an experience-driven RAG problem rather than a purely generative or purely heuristic task. Second, it introduces a reward-aware retrieval-and-reasoning architecture tailored to the structure of historical disruption cases. Third, it provides an empirical analysis on a benchmark with 21 event types, five severity levels, and repeated-query stability testing. Fourth, it shows that the accuracy gains of reward-aware reasoning can be obtained while keeping end-to-end latency in the single-digit millisecond range.

\section{Related Work}
\label{sec:related_work}

\subsection{Case-Based Reasoning for Operational Decisions}

CBR remains one of the most natural paradigms for operational decision support because it explicitly models experiential knowledge. Foundational work by Aamodt and Plaza \cite{aamodt1994cbr} and Watson \cite{watson1999cbr} established the retrieve--reuse--revise--retain cycle that still underpins modern memory-based systems. In supply-chain management, CBR has been used to coordinate distributed actors under uncertainty, as illustrated by the MACE-SCM framework \cite{kwon2007mace}. Fang and Wong \cite{fang2010hybrid} extended the idea to buyer--seller negotiations by combining CBR with multi-agent negotiation support, while Giannakis and Louis \cite{giannakis2011risk} used intelligent agent design for supply-chain risk management. These studies demonstrate that operational memory is valuable, but they also expose a recurring bottleneck: retrieval is usually dominated by hand-tuned features or symbolic similarity rules that do not scale well to heterogeneous event descriptions.

\subsection{Dense Retrieval and Retrieval-Augmented Generation}

RAG reframes memory use in language systems by combining a retriever with a generator \cite{lewis2020rag}. The retriever grounds the response in external evidence, while the generator synthesizes or ranks candidate outputs conditioned on the retrieved context. Dense Passage Retrieval (DPR) showed how dual-encoder representations support semantic matching beyond sparse lexical overlap \cite{karpukhin2020dpr}. Sentence-BERT demonstrated efficient sentence embeddings for semantic search and clustering \cite{reimers2019sbert}, and ColBERT introduced late interaction to preserve token-level matching while keeping retrieval scalable \cite{khattab2020colbert}. ColBERTv2 later improved the effectiveness--efficiency trade-off through compression and denoised supervision \cite{santhanam2022colbertv2}. These methods are directly relevant to logistics disruption management, where a query such as \textit{road closed near urban bridge under stress operations} should retrieve analogous experience even if historical cases are phrased differently.

Beyond text-only retrieval, recent work has begun to structure RAG around explicit relational knowledge. Knowledge-graph-guided RAG expands seed evidence through graph structure to improve contextual coherence \cite{zhu2025kg2rag}. This direction is important for logistics because disruption decisions are rarely determined by isolated facts; they depend on relations among event type, severity, scenario scale, available actions, and prior outcomes. Our approach is consistent with this trend, although it uses a case memory rather than a standalone knowledge graph as the primary non-parametric store.

\subsection{AI for Logistics, Resilience, and Disruption Management}

The broader logistics and supply-chain literature has documented a growing role for AI and machine learning in forecasting, routing, scheduling, monitoring, and resilience analysis \cite{toorajipour2021ai,akbari2021ml}. At the same time, disruption-focused studies emphasize recovery and ripple effects, showing that local failures propagate across upstream and downstream decisions \cite{ivanov2017recovery,dolgui2018ripple}. These findings motivate decision-support systems that can react at execution time, not only at planning time.

Operations research remains essential in this space. Vehicle-routing references provide the optimization backbone \cite{toth2014vrp,pillac2013dvrp}, and modern solvers such as CP-SAT in OR-Tools provide powerful exact and hybrid search capabilities \cite{perron2023cpsat}. Yet the disruption layer often requires semantic interpretation before an optimization model can even be instantiated. Our work therefore occupies the interface between AI and OR: it does not replace planning models, but supplies a fast and experience-grounded action prior that can trigger or guide subsequent optimization.

\subsection{Reasoning and Lightweight Deployment}

Once retrieval quality becomes sufficiently strong, the next issue is how to convert retrieved evidence into a robust action. Chain-of-thought prompting and tool-use studies suggest that language models can perform structured reasoning over external context \cite{wei2022cot,schick2023toolformer}. For operations, however, the model must also be deployable under practical latency and hardware constraints. Work on open foundation models and quantized adaptation shows that compact reasoning stacks can be used without the full cost profile of frontier-scale deployment \cite{touvron2023llama,dettmers2023qlora,lin2024awq}. This makes experience-driven RAG particularly appealing for logistics: the system can reserve expensive optimization for the few cases that truly need it, while still exploiting language-based reasoning for most disruptions.

The literature therefore suggests a clear gap. CBR provides explainable experiential memory, dense retrieval improves semantic access to that memory, and lightweight language reasoning can turn retrieved evidence into an action. What remains underexplored is a unified design for logistics disruption management that explicitly evaluates decision accuracy, retrieval quality, stability, and latency together. This paper addresses that gap.

\section{System Architecture}
\label{sec:architecture}

\begin{figure*}[t]
\centering
\includegraphics[width=\linewidth,height=2in,keepaspectratio]{\casefig}
\caption{Case distribution across 21 event types and four operational scenarios in the 1,000-case library.}
\label{fig:case_distribution}
\end{figure*}

The proposed system organizes logistics decision support around an explicit experience memory and a staged reasoning pipeline. Rather than directly prompting a model with the incoming disruption, we first identify the most relevant prior experiences, then transform those experiences into a grounded action recommendation. This architecture is motivated by three operational requirements: millisecond-level responsiveness, traceability to historical precedents, and robustness to scenario shifts.

\subsection{Case Memory}

Each case stores a disruption event, contextual descriptors, the historical action, and the observed outcome. The case library comprises 1,000 cases, 21 event types, four scenarios, and five admissible actions: \textit{reroute}, \textit{ignore}, \textit{adjust\_capacity}, \textit{reassign\_order}, and \textit{delay\_tolerant}. The outcome field is binary, with 767 successful cases and 233 failed cases, yielding a 76.7\% success rate. This is important because the system should not merely copy common actions; it should prefer actions that worked under comparable circumstances.

The memory exhibits two complementary distributional properties. First, scenarios are evenly represented with 250 cases each, which limits bias toward a particular operating scale. Second, event frequencies are highly unequal. \textit{Vehicle\_breakdown} contributes 361 cases, or 36.1\% of the library, while the next largest categories are far smaller. This long-tail structure is realistic for logistics operations and creates a difficult retrieval problem: the system must handle dominant routine incidents without collapsing rare-event discrimination.

\subsection{Pipeline Overview}

The architecture contains four stages.

\textit{Stage 1: Query normalization.} Incoming disruptions are converted into a structured representation containing event semantics, severity, and operational scenario. The benchmark covers 420 such queries, corresponding to 21 event types, four scenarios, and five severity levels.

\textit{Stage 2: Experience retrieval.} The normalized query is matched against the case memory. The baseline uses a manually designed five-dimensional score. V1 and V2 replace purely heuristic matching with semantically richer retrieval while retaining operational factors such as severity and action compatibility.

\textit{Stage 3: Action reasoning.} Retrieved cases are interpreted as evidence for a discrete action choice. V1 predicts the action from the retrieved support set. V2 additionally evaluates whether the candidate action is consistent with successful historical outcomes.

\textit{Stage 4: Decision output.} The system returns the top action and the top-ranked supporting cases. This makes the output inspectable by dispatchers or by downstream optimization modules.

\subsection{Why Experience-Driven RAG Fits Disruption Management}

Disruption management differs from open-domain question answering in two ways. First, the output space is constrained: the system selects one of five actions rather than generating free text. Second, the memory items are operational episodes, not documents. This means the retriever must maximize utility rather than textual similarity alone. For example, two cases may describe different event labels but still imply the same response if they create comparable route infeasibility or service risk.

The architecture addresses this by treating retrieval as action-oriented evidence gathering. Dense semantic relevance identifies analogous contexts, but the downstream scoring stage incorporates structured logistics features. In effect, the system uses RAG not to produce long-form explanations, but to ground a compact operational decision in the most relevant prior experience.

\subsection{Baseline, V1, and V2 as Architectural Variants}

The three evaluated methods can be understood as architectural ablations.

The \textit{Baseline} uses only a five-dimensional score. It is computationally cheap and responds in 0.495\,ms on average, but it has no explicit mechanism for semantic abstraction beyond the engineered feature space.

\textit{V1} adds a reasoning layer on top of retrieved cases. This produces an action hypothesis, but because the final choice is not filtered by explicit historical success evidence, it can over-trust retrieved narratives that are similar yet not operationally validated.

\textit{V2} adds reward-aware reasoning. Instead of asking only whether a case looks similar, it asks whether the proposed action is supported by successful outcomes under comparable conditions. This extra reasoning stage raises the average decision time to 8.415\,ms, but it also produces the highest action accuracy and NDCG@5.

\subsection{Support for Human and Algorithmic Supervision}

An operational decision-support system must justify its recommendation to both people and machines. The case-memory design serves both audiences. For human dispatchers, the retrieved cases show why a recommendation was made. For downstream solvers, the decision can act as a warm-start or rule-based prior. If a routing or scheduling module is later invoked, the recommended action narrows the search space. This is especially useful in stress scenarios, where full replanning may be too costly to run for every event.

\begin{table}[t]
\caption{Benchmark and Case-Memory Profile}
\label{tab:profile}
\centering
\small
\begin{tabular}{ll}
\toprule
\textbf{Item} & \textbf{Value} \\
\midrule
Case library size & 1,000 cases \\
Event types & 21 \\
Operational scenarios & 4 (small, medium, large, stress) \\
Candidate actions & 5 \\
Outcome success rate & 76.7\% \\
Benchmark queries & 420 \\
Severity levels & 5 \\
Stability repeats/query & 10 \\
\bottomrule
\end{tabular}
\end{table}

Table~\ref{tab:profile} summarizes the operational scope of the benchmark. The benchmark is deliberately compact enough for repeated evaluation yet diverse enough to stress both semantic matching and action consistency.

\begin{figure*}[t]
\centering
\includegraphics[width=\linewidth,height=2in,keepaspectratio]{\compfig}
\caption{Comparison of Baseline, V1, and V2 on action accuracy, NDCG@5, and average decision time.}
\label{fig:comparison}
\end{figure*}

\section{Methodology}
\label{sec:method}

The decision process can be formalized as retrieval followed by reward-aware action selection. Let $q$ denote an incoming disruption query and let $\mathcal{C}$ be the case memory. Each case $c \in \mathcal{C}$ contains an event description, context tuple, historical action $a_c$, and outcome $o_c$.

\subsection{Action-Oriented Retrieval}

The baseline score is designed to capture operational relevance rather than plain text similarity. We write the composite relevance as
\begin{align*}
s(q,c) &= \lambda_1 s_{\mathrm{sem}}(q,c) + \lambda_2 s_{\mathrm{sev}}(q,c) \\
&\quad + \lambda_3 s_{\mathrm{scn}}(q,c) + \lambda_4 s_{\mathrm{act}}(q,c)
+ \lambda_5 s_{\mathrm{fea}}(q,c),
\end{align*}
where $s_{\mathrm{sem}}$ is semantic similarity, $s_{\mathrm{sev}}$ measures severity alignment, $s_{\mathrm{scn}}$ captures scenario compatibility, $s_{\mathrm{act}}$ reflects action compatibility, and $s_{\mathrm{fea}}$ represents feasibility cues derived from structured context.

In the baseline, these terms are implemented through manually designed matching rules. In V1 and V2, the semantic term is upgraded through dense retrieval so that the candidate set better reflects latent similarity between disruption situations. The retrieved top-$k$ cases define an evidence set $\mathcal{R}_k(q)$.

\subsection{Action Aggregation}

For each action $a$ in the candidate action set $\mathcal{A}$, the retrieved evidence can be aggregated as
\begin{align*}
\phi(a \mid q) = \sum_{c \in \mathcal{R}_k(q)} \mathbb{I}[a_c = a] \cdot s(q,c),
\end{align*}
where $\mathbb{I}[\cdot]$ is the indicator function. The baseline effectively selects the action with the highest aggregated support in the retrieved neighborhood. V1 augments this step with a reasoning module that interprets the retrieved cases rather than simply counting or summing support.

\subsection{Reward-Aware Reranking}

The key addition in V2 is a reward-aware decision layer. Historical outcomes let us distinguish actions that were merely common from actions that were successful. We therefore define a reward-adjusted score
\begin{align*}
\psi(a \mid q) = \phi(a \mid q) + \beta \cdot r(a,q,\mathcal{R}_k(q)),
\end{align*}
where $r(\cdot)$ estimates how well action $a$ is supported by successful retrieved cases under the current disruption context. Conceptually, this term favors actions whose historical support is both similar and effective.

This design is important in imbalanced libraries. When a frequent event type dominates the memory, a pure nearest-neighbor strategy can overfit to majority behavior. Reward-aware reranking counteracts this by emphasizing outcome quality, not only exposure frequency.

\subsection{Decision Time and Stability}

The evaluation uses two classes of metrics. The first class measures decision quality through action accuracy and NDCG@5. The second measures operational reliability through average decision time and repeated-query stability. Stability is estimated by running each query 10 times and computing Top-1 action consistency and Top-5 exact consistency. These metrics matter because dispatch environments favor consistent recommendations; an unstable retriever can undermine user trust even when average accuracy is high.

\subsection{Interpretation of the Three Variants}

Methodologically, the benchmark isolates the value of successive reasoning layers. The baseline tests whether carefully designed structured scoring is sufficient. V1 tests whether a retrieval-conditioned action predictor alone improves decisions. V2 tests whether adding reward awareness creates a better bias toward historically successful interventions. The resulting comparison shows that stronger reasoning is beneficial only when it remains grounded in outcome-aware experience.

\section{Experiments}
\label{sec:experiments}

\subsection{Experimental Setting}

The benchmark contains 420 disruption queries derived from 21 event types, four scenarios, and five severity levels. This ensures that every event type is evaluated across multiple operating scales, from small to stress conditions. The case memory contains 1,000 cases with the same event vocabulary and action space, enabling direct comparison between retrieved evidence and ground-truth decisions.

We evaluate three methods:
\begin{itemize}
\item \textbf{Baseline}: five-dimensional scoring with no reward-aware reasoning.
\item \textbf{V1}: retrieval plus action prediction.
\item \textbf{V2}: retrieval plus action prediction and reward-aware reranking.
\end{itemize}

Action Accuracy measures whether the top action matches the reference action. NDCG@5 measures the ranking quality of the top-5 retrieved support set. Average decision time reports end-to-end latency in milliseconds. Stability is measured by repeating each query 10 times.

\subsection{Main Results}

Table~\ref{tab:main_results} reports the core comparison. V2 attains the highest action quality and retrieval quality, reaching 99.762\% action accuracy and 99.685\% NDCG@5. Compared with the baseline, this corresponds to gains of 4.524 percentage points in action accuracy and 2.910 percentage points in NDCG@5. Compared with V1, the gains are even larger: 7.857 points in accuracy and 4.951 points in NDCG@5.

\begin{table}[t]
\caption{Overall Performance on 420 Benchmark Queries}
\label{tab:main_results}
\centering
\small
\begin{tabular}{lccc}
\toprule
\textbf{Method} & \textbf{Accuracy} & \textbf{NDCG@5} & \textbf{Time (ms)} \\
\midrule
Baseline & 95.238\% & 96.775\% & 0.495 \\
V1 & 91.905\% & 94.734\% & 2.271 \\
V2 (ours) & \textbf{99.762\%} & \textbf{99.685\%} & 8.415 \\
\bottomrule
\end{tabular}
\end{table}

The latency trade-off is visible but operationally acceptable. The baseline is the fastest at 0.495\,ms, and V2 is the slowest at 8.415\,ms. Yet V2 remains comfortably below 10\,ms, which is fast enough for interactive dispatch support and for integration into larger planning loops. Figure~\ref{fig:comparison} visualizes this trade-off: the extra computation buys a very large gain in decision quality.

One notable result is that V1 performs worse than the baseline on both accuracy and NDCG@5. This suggests that adding an unconstrained reasoning layer is not sufficient by itself. Without outcome-aware reward shaping, the action predictor can amplify errors from imperfect retrieval or over-interpret marginal evidence. V2 corrects this by forcing the final decision to remain aligned with historically successful cases.

\subsection{Retrieval Stability}

Figure~\ref{fig:stability} and Table~\ref{tab:stability} present repeated-query stability. V2 achieves 99.905\% Top-1 action consistency, compared with 99.143\% for V1 and 98.500\% for the baseline. The gain over the baseline is 1.405 percentage points. V2 also records the best Top-5 exact consistency at 16.881\%, exceeding the baseline by 1.738 points.

\begin{figure}[t]
\centering
\includegraphics[width=\linewidth,height=2in,keepaspectratio]{\stabfig}
\caption{Retrieval stability under 10 repeated runs per query.}
\label{fig:stability}
\end{figure}

\begin{table}[t]
\caption{Retrieval Stability over 10 Repeated Runs per Query}
\label{tab:stability}
\centering
\small
\begin{tabular}{lcc}
\toprule
\textbf{Method} & \textbf{Top-1 Consistency} & \textbf{Top-5 Exact} \\
\midrule
Baseline & 98.500\% & 15.143\% \\
V1 & 99.143\% & 15.929\% \\
V2 (ours) & \textbf{99.905\%} & \textbf{16.881\%} \\
\bottomrule
\end{tabular}
\end{table}

The contrast between the two stability metrics is informative. Top-1 consistency is almost perfect, while Top-5 exact consistency is low for all methods. This means the best action is highly stable even when the tail of the support set changes order. In practice, that is a favorable failure mode: dispatchers care primarily about whether the top recommendation is reproducible, not whether the fourth and fifth evidence cases swap positions.

\subsection{Case-Library Effects}

The case distribution in Figure~\ref{fig:case_distribution} helps explain why reward-aware retrieval matters. The library is balanced across scenarios but not across events. The dominance of \textit{vehicle\_breakdown} creates a strong prior toward routing and capacity interventions, while rare events such as natural disasters or weather-delay variants require sharper semantic discrimination. A purely frequency-driven retriever could therefore drift toward common operational responses. V2 mitigates this risk by scoring evidence through both similarity and historical success.

Action frequencies are also uneven: \textit{reroute} appears 350 times, \textit{delay\_tolerant} 230 times, \textit{adjust\_capacity} 186 times, \textit{reassign\_order} 140 times, and \textit{ignore} 94 times. This skew is typical of logistics execution, where active interventions dominate passive ones. The near-perfect accuracy of V2 indicates that it is not merely following marginal action frequencies; otherwise its performance on minority actions would degrade more visibly.

\subsection{Implications for Operations}

From an operational perspective, the results support three conclusions. First, structured heuristics remain competitive when latency is the only concern. Second, retrieval-conditioned reasoning is useful only if it is grounded by experience quality rather than narrative plausibility. Third, a sub-10\,ms reward-aware RAG stack can serve as an execution-time decision layer before heavier optimization is invoked. This makes V2 suitable as a triage mechanism for disruption management systems in logistics control towers.

\section{Conclusion}
\label{sec:conclusion}

This paper presented an experience-driven retrieval-augmented generation framework for logistics disruption management. The proposed design combines a structured case memory, semantically informed retrieval, and reward-aware reasoning over historical outcomes. On a benchmark of 420 queries drawn from 21 event types and four operating scenarios, V2 achieved 99.762\% action accuracy, 99.685\% NDCG@5, and 99.905\% Top-1 action consistency, while maintaining an average decision time of 8.415\,ms.

The empirical pattern is clear. Hand-crafted retrieval is fast but limited, action prediction alone is insufficient, and reward-aware reasoning is the component that turns retrieved experience into consistently correct decisions. The results show that experience-driven RAG is a practical design for logistics disruption support when the output space is discrete, historical memory is available, and traceability matters.

Future work should extend the memory with temporal dependencies and multi-event cascades, connect the action layer to downstream optimization loops, and test closed-loop adaptation as new cases are retained over time. A second direction is to enrich the case memory with graph structure so that interdependent disruptions can be retrieved jointly rather than case by case.

\bibliographystyle{IEEEtran}
\bibliography{main}

\end{document}
